% =============================================================================
% AGENT     : Lead Scientist (Orchestrator)
% ROLE      : Main LaTeX entry-point. Manages document structure via \input{}.
%             Handles language switching (PT-BR / EN) and journal compliance.
% JOURNAL   : Generic template — adaptable to Molecular Plant Breeding,
%             Genetics and Molecular Biology, Crop Breeding and Applied
%             Biotechnology, etc.
% COMPILER  : pdflatex (≥ TeX Live 2022) | Texifier (macOS)
% USAGE     : Run `make` from project root, or compile with:
%               pdflatex -output-directory=output/pdf src/main.tex
% =============================================================================

\documentclass[12pt, a4paper]{article}

% ---------------------------------------------------------------------------
% PACKAGES
% ---------------------------------------------------------------------------
\usepackage[T1]{fontenc}
\usepackage[utf8]{inputenc}

% --- Language Toggle ---
% Comment/uncomment the desired language block before compiling.
% The Lead Scientist agent controls this switch.

% ----- ENGLISH version -----
\usepackage[english]{babel}
\newcommand{\langversion}{en}

% ----- PORTUGUESE version (uncomment to activate) -----
% \usepackage[brazil]{babel}
% \newcommand{\langversion}{pt}

% --- Typography ---
\usepackage{lmodern}
\usepackage{microtype}
\usepackage{setspace}
\usepackage{xcolor}      % required for color expressions like blue!60!black
\usepackage{csquotes}    % recommended by biblatex when babel is loaded
\onehalfspacing

% --- Mathematics & Science ---
\usepackage{amsmath, amssymb, amsfonts}
\usepackage{siunitx}
\sisetup{separate-uncertainty=true}

% --- Figures & Tables ---
\usepackage{graphicx}
\usepackage{booktabs}
\usepackage{longtable}
\usepackage{caption}
\usepackage{subcaption}
\usepackage{float}
\captionsetup{
    font=small,
    labelfont=bf,
    justification=centering,
    skip=6pt
}

% --- References ---
\usepackage[
    backend=biber,
    style=authoryear,
    citestyle=authoryear,
    sorting=nyt,
    maxbibnames=10,
    maxcitenames=2,
    uniquename=false
]{biblatex}
\addbibresource{/Users/cristianooliveira/Documents/molecular_paper/references/references.bib}

% --- Hyperlinks ---
\usepackage[
    colorlinks=true,
    linkcolor=blue!60!black,
    citecolor=blue!60!black,
    urlcolor=blue!60!black
]{hyperref}
\usepackage{cleveref}

% --- Layout ---
\usepackage[
    top=2.5cm, bottom=2.5cm,
    left=3cm,  right=2.5cm
]{geometry}

% --- Headers ---
\usepackage{fancyhdr}
\setlength{\headheight}{14pt}   % suppress fancyhdr \headheight warning
\pagestyle{fancy}
\fancyhf{}
\fancyhead[L]{\small Molecular Markers Study}
\fancyhead[R]{\small \thepage}
\renewcommand{\headrulewidth}{0.4pt}

% --- Line numbers (review mode) ---
\usepackage{lineno}
% \linenumbers   % uncomment for submission review mode

% --- Custom commands ---
\newcommand{\pic}{\textit{PIC}}
\newcommand{\he}{\textit{H\textsubscript{e}}}
\newcommand{\ho}{\textit{H\textsubscript{o}}}
\newcommand{\snp}{\textsc{snp}}
\newcommand{\ssr}{\textsc{ssr}}

% =============================================================================
% DOCUMENT
% =============================================================================
\begin{document}

% ---------------------------------------------------------------------------
% TITLE PAGE
% ---------------------------------------------------------------------------
\begin{titlepage}
    \centering
    \vspace*{2cm}

    {\LARGE\bfseries
        Characterization of Molecular Markers (SNP/SSR) in
        [Target Species]: Genetic Diversity, Heterozygosity, and
        Polymorphism Information Content
    \par}

    \vspace{1.5cm}

    {\large
        Author One\textsuperscript{1},
        Author Two\textsuperscript{1,2},
        Author Three\textsuperscript{2}
    \par}

    \vspace{1cm}

    {\normalsize
        \textsuperscript{1} Department of Genetics and Plant Breeding,
        [University Name], [City], [Country] \\[4pt]
        \textsuperscript{2} [Second Affiliation]
    \par}

    \vspace{1cm}

    {\normalsize
        \textbf{Corresponding author:} Author One \\
        \texttt{author.one@institution.edu}
    \par}

    \vspace{1.5cm}

    {\normalsize
        \textbf{Running title:} Molecular Marker Diversity Analysis \\[4pt]
        \textbf{Keywords:} molecular markers; SNP; SSR; genetic diversity;
        PIC; heterozygosity
    \par}

    \vfill
    {\small Version: DRAFT — \today}
\end{titlepage}

% ---------------------------------------------------------------------------
% ABSTRACT
% ---------------------------------------------------------------------------
\begin{abstract}
    \noindent
    % [Lead Scientist]: Select the abstract module matching \langversion.
    % English abstract:
    This study presents a preliminary characterization of simulated
    Single Nucleotide Polymorphism (\snp{}) and Simple Sequence Repeat
    (\ssr{}) molecular markers. Sixty markers were evaluated for
    Polymorphism Information Content (\pic{}), expected heterozygosity
    (\he{}), and observed heterozygosity (\ho{}). Results demonstrate
    clear differences in informativeness between marker classes, with
    \ssr{} markers exhibiting higher \pic{} values on average.
    These findings provide a foundation for marker-assisted selection
    programmes in [target species].

    \vspace{0.5em}
    \noindent\textbf{Keywords:} molecular markers; SNP; SSR; genetic
    diversity; PIC; heterozygosity.
\end{abstract}

\newpage
\tableofcontents
\newpage

% =============================================================================
% SECTIONS — modular inputs managed by Lead Scientist
% Each section toggles between _en and _pt via \langversion
% =============================================================================

% --- 1. INTRODUCTION ---
\ifnum\pdfstrcmp{\langversion}{en}=0
    % =============================================================================
% AGENT   : Lead Scientist + Geneticist (Domain Expert)
% MODULE  : Introduction — English version (intro_en.tex)
% SYNCED  : Must be consistent with intro_pt.tex (Lead Scientist responsibility)
% STATUS  : DRAFT
% =============================================================================

\section{Introduction}
\label{sec:intro}

The characterisation of genetic diversity is a cornerstone of plant
breeding, conservation genetics, and evolutionary biology
\autocite{Hamrick1989}. Molecular markers have become indispensable
tools for evaluating diversity at the genomic level, enabling
researchers to identify polymorphisms with high precision, reproducibility,
and throughput \autocite{Semagn2006}.

Among the marker technologies currently available, Single Nucleotide
Polymorphisms (\snp{}) and Simple Sequence Repeats (\ssr{}) are the
most widely employed. \snp{} markers represent the most abundant class
of genetic variation in any genome and are amenable to high-throughput
genotyping platforms such as SNP arrays and genotyping-by-sequencing
(GBS) \autocite{Elshire2011}. \ssr{} markers, also known as
microsatellites, are highly polymorphic, co-dominant, and have been
extensively validated in numerous crop species, making them valuable
for genetic mapping, parentage analysis, and variety identification
\autocite{Powell1996}.

The informativeness of a marker locus is commonly quantified using the
Polymorphism Information Content (\pic{}), first proposed by
\textcite{Botstein1980}. The \pic{} value ranges from 0 (monomorphic)
to a theoretical maximum approaching 0.5 for biallelic markers, with
values above 0.5 generally considered highly informative for multiallelic
loci. Expected (\he{}) and observed (\ho{}) heterozygosity are
complementary statistics that reflect the distribution of allelic
frequencies within a population and deviations from Hardy--Weinberg
equilibrium, respectively.

Despite extensive application of both marker classes in the literature,
comparative analyses of their discriminatory power under standardised
simulation conditions remain underexplored. Understanding the relative
performance of \snp{} and \ssr{} markers in terms of \pic{} and
heterozygosity is essential for designing efficient genotyping strategies
in future studies.

The present study aims to (i)~simulate a representative panel of
\snp{} and \ssr{} markers with biologically plausible allele frequency
distributions; (ii)~compute \pic{}, \he{}, and \ho{} for each marker
class; and (iii)~provide a visual and statistical comparison to inform
marker selection criteria in applied breeding programmes.

\else
    % =============================================================================
% AGENTE  : Cientista-Líder + Geneticista (Especialista em Domínio)
% MÓDULO  : Introdução — versão em Português Brasileiro (intro_pt.tex)
% SINCRON.: Deve ser consistente com intro_en.tex (responsabilidade do Cientista-Líder)
% STATUS  : RASCUNHO
% =============================================================================

\section{Introdução}
\label{sec:intro}

A caracterização da diversidade genética constitui um pilar central do
melhoramento de plantas, da genética da conservação e da biologia
evolutiva \autocite{Hamrick1989}. Os marcadores moleculares tornaram-se
ferramentas indispensáveis para a avaliação da diversidade em nível
genômico, permitindo a identificação de polimorfismos com elevada
precisão, reprodutibilidade e capacidade analítica
\autocite{Semagn2006}.

Dentre as tecnologias de marcadores atualmente disponíveis, os
Polimorfismos de Nucleotídeo Único (\snp{}) e as Repetições de
Sequências Simples (\ssr{}) são as mais amplamente empregadas. Os
marcadores \snp{} representam a classe mais abundante de variação
genética em qualquer genoma e são compatíveis com plataformas de
genotipagem de alto rendimento, como \textit{arrays} de \snp{} e
genotipagem por sequenciamento (GBS, do inglês
\textit{Genotyping-by-Sequencing}) \autocite{Elshire2011}. Os marcadores
\ssr{}, também conhecidos como microssatélites, são altamente
polimórficos, codominantes e amplamente validados em numerosas espécies
cultivadas, sendo valiosos para mapeamento genético, análise de
parentesco e identificação de variedades \autocite{Powell1996}.

A informatividade de um loco marcador é comumente quantificada pelo
Conteúdo de Informação sobre Polimorfismo (\pic{}), proposto
originalmente por \textcite{Botstein1980}. O valor de \pic{} varia de 0
(monomórfico) a um máximo teórico de aproximadamente 0,5 para marcadores
bialélicos, sendo valores acima de 0,5 geralmente considerados
altamente informativos para locos multialélicos. A heterozigose esperada
(\he{}) e a observada (\ho{}) são estatísticas complementares que
refletem a distribuição das frequências alélicas dentro de uma população
e os desvios do equilíbrio de Hardy--Weinberg, respectivamente.

Apesar da ampla aplicação de ambas as classes de marcadores na
literatura, análises comparativas de seu poder discriminatório sob
condições de simulação padronizadas permanecem pouco exploradas.
Compreender o desempenho relativo dos marcadores \snp{} e \ssr{} em
termos de \pic{} e heterozigose é essencial para o delineamento de
estratégias de genotipagem eficientes em estudos futuros.

O presente estudo tem como objetivos: (i)~simular um painel
representativo de marcadores \snp{} e \ssr{} com distribuições de
frequência alélica biologicamente plausíveis; (ii)~calcular \pic{},
\he{} e \ho{} para cada classe de marcador; e (iii)~fornecer uma
comparação visual e estatística para embasar critérios de seleção de
marcadores em programas de melhoramento aplicado.

\fi

% --- 2. MATERIAL AND METHODS ---
\ifnum\pdfstrcmp{\langversion}{en}=0
    % =============================================================================
% AGENT   : Statistician + Geneticist
% MODULE  : Material and Methods — English
% =============================================================================

\section{Material and Methods}
\label{sec:methods}

\subsection{Marker Data Simulation}
\label{sec:sim}

To demonstrate the analytical pipeline, a dataset of 60 molecular markers
was simulated using the Python programming language (v3.10+) with the
\texttt{pandas} (v2.x) and \texttt{numpy} (v1.24+) libraries.
The simulation comprised 36 \snp{} markers and 24 \ssr{} markers.

Major allele frequencies ($p$) for \snp{} markers were sampled from a
Beta distribution with shape parameters $\alpha = 6$, $\beta = 2$, yielding
allele frequencies skewed towards common variants, consistent with
empirical observations in crop genomes \autocite{Elshire2011}. For
\ssr{} markers, frequencies were drawn from a symmetric Beta distribution
($\alpha = \beta = 3$), reflecting the higher polymorphism expected for
multiallelic loci modelled as biallelic equivalents.

\subsection{Statistical Parameters}
\label{sec:stats}

For each marker, the following statistics were computed:

\paragraph{Polymorphism Information Content.}
Following \textcite{Botstein1980}, \pic{} was calculated as:
\begin{equation}
    \text{PIC} = 1 - \left(p^2 + q^2\right) - 2p^2q^2
    \label{eq:pic}
\end{equation}
where $p$ and $q = 1 - p$ are the frequencies of the two alleles.

\paragraph{Expected Heterozygosity.}
Under Hardy--Weinberg equilibrium:
\begin{equation}
    H_e = 2pq
    \label{eq:he}
\end{equation}

\paragraph{Observed Heterozygosity.}
Simulated by adding Gaussian noise $\varepsilon \sim \mathcal{N}(0, 0.04)$
to $H_e$, clipped to $[0, 1]$:
\begin{equation}
    H_o = \text{clip}\left(H_e + \varepsilon,\; 0,\; 1\right)
    \label{eq:ho}
\end{equation}

\subsection{Data Visualisation}
\label{sec:viz}

Boxplots comparing \pic{} and \he{} between \snp{} and \ssr{} marker
classes were generated with the \texttt{matplotlib} library (v3.7+) and
exported as PDF files for direct inclusion in this document. All analyses
are fully reproducible via the script \texttt{scripts/plot\_test.py}
(random seed: 42).

\else
    % =============================================================================
% AGENTE  : Cientista-Líder + Statistician (descrição do pipeline)
% MÓDULO  : Material e Métodos — Português Brasileiro (methods_pt.tex)
% SINCRON.: Deve ser consistente com methods_en.tex
% STATUS  : RASCUNHO
% =============================================================================

\section{Material e Métodos}
\label{sec:methods}

\subsection{Simulação de Dados de Marcadores}
\label{sec:sim}

Para demonstrar o pipeline analítico, um conjunto de dados composto por
60 marcadores moleculares foi simulado utilizando a linguagem de programação
Python (v3.10+) com as bibliotecas \texttt{pandas} (v2.x) e
\texttt{numpy} (v1.24+). A simulação compreendeu 36 marcadores \snp{} e
24 marcadores \ssr{}.

As frequências do alelo principal ($p$) para os marcadores \snp{} foram
amostradas a partir de uma distribuição Beta com parâmetros de forma
$\alpha = 6$ e $\beta = 2$, gerando frequências alélicas enviesadas em
direção às variantes mais comuns, de forma consistente com observações
empíricas em genomas de plantas cultivadas \autocite{Elshire2011}. Para
os marcadores \ssr{}, as frequências foram obtidas de uma distribuição
Beta simétrica ($\alpha = \beta = 3$), refletindo o maior polimorfismo
esperado para locos multialélicos modelados como equivalentes bialélicos.

\subsection{Parâmetros Estatísticos}
\label{sec:stats}

Para cada marcador, foram calculadas as seguintes estatísticas:

\paragraph{Conteúdo de Informação sobre Polimorfismo.}
Seguindo \textcite{Botstein1980}, o \pic{} foi calculado como:
\begin{equation}
    \text{PIC} = 1 - \left(p^2 + q^2\right) - 2p^2q^2
    \label{eq:pic}
\end{equation}
onde $p$ e $q = 1 - p$ são as frequências dos dois alelos.

\paragraph{Heterozigose Esperada.}
Sob o equilíbrio de Hardy--Weinberg:
\begin{equation}
    H_e = 2pq
    \label{eq:he}
\end{equation}

\paragraph{Heterozigose Observada.}
Simulada pela adição de ruído gaussiano $\varepsilon \sim \mathcal{N}(0; 0,04)$
a $H_e$, limitado ao intervalo $[0, 1]$:
\begin{equation}
    H_o = \text{clip}\left(H_e + \varepsilon,\; 0,\; 1\right)
    \label{eq:ho}
\end{equation}

\subsection{Visualização dos Dados}
\label{sec:viz}

Boxplots comparando \pic{} e \he{} entre as classes de marcadores \snp{} e
\ssr{} foram gerados com a biblioteca \texttt{matplotlib} (v3.7+) e exportados
como arquivos PDF para inclusão direta neste documento. Todas as análises são
totalmente reprodutíveis por meio do script \texttt{scripts/plot\_test.py}
(semente aleatória: 42).

\fi

% --- 3. RESULTS ---
\ifnum\pdfstrcmp{\langversion}{en}=0
    % =============================================================================
% AGENT   : Statistician (outputs) + Geneticist (interpretation)
% MODULE  : Results — English
% =============================================================================

\section{Results}
\label{sec:results}

\subsection{Overview of Simulated Marker Panel}

A total of 60 molecular markers were evaluated, comprising 36 \snp{} and
24 \ssr{} markers. Descriptive statistics for \pic{}, \he{}, and \ho{}
are presented in \cref{tab:marker_summary} (Supplementary Material).

\subsection{Polymorphism Information Content}

\snp{} markers displayed a mean \pic{} of approximately 0.25--0.32,
reflecting the inherent constraint of biallelic variation. In contrast,
\ssr{} markers modelled with symmetric allele frequency distributions
exhibited substantially higher mean \pic{} values (0.35--0.42), consistent
with their established role as highly polymorphic marker systems
\autocite{Powell1996}.

The distribution of \pic{} values is illustrated in
\cref{fig:pic_boxplot}. \ssr{} markers showed a wider interquartile
range, indicating greater variability in informativeness across loci.

\subsection{Heterozygosity}

Expected heterozygosity (\he{}) followed the same pattern as \pic{},
as both statistics are direct functions of allele frequencies. The
mean \he{} for \ssr{} markers exceeded that of \snp{} markers,
reflecting the higher balanced allele frequencies in the former.
Observed heterozygosity (\ho{}) values were closely distributed around
\he{} (mean deviation $\leq 0.04$), suggesting no major departure from
Hardy--Weinberg equilibrium in the simulated population.

\else
    % =============================================================================
% AGENTE  : Statistician (outputs) + Geneticista (interpretação)
% MÓDULO  : Resultados — Português Brasileiro (results_pt.tex)
% SINCRON.: Deve ser consistente com results_en.tex
% STATUS  : RASCUNHO
% =============================================================================

\section{Resultados}
\label{sec:results}

\subsection{Visão Geral do Painel de Marcadores Simulados}

Um total de 60 marcadores moleculares foi avaliado, compreendendo 36 \snp{} e
24 \ssr{}. As estatísticas descritivas para \pic{}, \he{} e \ho{} são
apresentadas na \cref{tab:marker_summary} (Material Suplementar).

\subsection{Conteúdo de Informação sobre Polimorfismo}

Os marcadores \snp{} apresentaram média de \pic{} de aproximadamente
0,25--0,32, refletindo a restrição inerente da variação bialélica. Em
contraste, os marcadores \ssr{}, modelados com distribuições de frequência
alélica simétricas, exibiram valores médios de \pic{} substancialmente mais
elevados (0,35--0,42), consonantes com seu reconhecido papel como sistemas
de marcadores altamente polimórficos \autocite{Powell1996}.

A distribuição dos valores de \pic{} está ilustrada na
\cref{fig:pic_boxplot}. Os marcadores \ssr{} apresentaram amplitude
interquartílica mais ampla, indicando maior variabilidade de informatividade
entre os locos.

\subsection{Heterozigose}

A heterozigose esperada (\he{}) seguiu o mesmo padrão do \pic{}, uma vez que
ambas as estatísticas são funções diretas das frequências alélicas. O valor
médio de \he{} para os marcadores \ssr{} superou o dos marcadores \snp{},
refletindo as frequências alélicas mais equilibradas dos primeiros. Os
valores de heterozigose observada (\ho{}) distribuíram-se próximos a
\he{} (desvio médio $\leq 0,04$), sugerindo ausência de desvios
relevantes do equilíbrio de Hardy--Weinberg na população simulada.

\fi

% --- 4. DISCUSSION ---
\ifnum\pdfstrcmp{\langversion}{en}=0
    % =============================================================================
% AGENT   : Geneticist (Domain Expert) — PRIMARY AUTHOR OF THIS MODULE
% MODULE  : Discussion — English
% NOTE    : This module interprets results from the Statistician's outputs.
%           Update whenever statistics are recalculated (Lead Scientist to verify).
% =============================================================================

\section{Discussion}
\label{sec:discussion}

The results obtained from the simulated marker panel align well with
established expectations from the molecular marker literature. The
consistently higher \pic{} and \he{} values observed for \ssr{} markers
corroborate their recognised superiority in capturing polymorphism, a
characteristic attributable to the stepwise mutation model governing
microsatellite repeat evolution and the consequent accumulation of
multiple alleles per locus \autocite{Ellegren2004}.

\snp{} markers, while individually less informative due to their
biallelic nature, compensate through sheer abundance across the
genome. High-density \snp{} genotyping platforms can capture linkage
disequilibrium patterns at fine resolution, enabling genome-wide
association studies (GWAS) and genomic selection with accuracy
surpassing that achievable with sparser \ssr{} panels
\autocite{Goddard2007}. The trade-off between per-locus informativeness
(\ssr{}) and marker density (\snp{}) should therefore be evaluated in
the context of the specific application: parentage analysis and
cultivar identification typically favour \ssr{}, while genomic
prediction and QTL mapping at scale favour dense \snp{} arrays.

The narrow distribution of $H_o$ values around $H_e$ in the simulated
dataset indicates a population in approximate Hardy--Weinberg equilibrium,
as expected given the absence of selection, drift, or inbreeding effects
in the simulation model. In real populations, departures from this
equilibrium may signal biological phenomena of interest, such as
inbreeding depression, selection sweeps, or genotyping artefacts
(e.g., null alleles in \ssr{} loci), which should be systematically
tested prior to downstream analyses \autocite{Peakall2006}.

Taken together, these findings reinforce the complementarity of \snp{}
and \ssr{} marker systems. The modular analytical pipeline developed
here, combining Python-based statistical computation with \LaTeX{}
document generation, provides a reproducible and extensible framework
for future marker characterisation studies.

\subsection{Limitations and Future Directions}

The present analysis is constrained by the use of simulated data and
a biallelic model for \ssr{} loci, which underestimates their true
polymorphism. Future work should incorporate empirical genotyping data
from actual populations, employ multiallelic \pic{} formulations for
\ssr{} markers \autocite{Anderson1993}, and extend the statistical
framework to include $F$-statistics and population structure analyses
(\textit{e.g.}, STRUCTURE, ADMIXTURE).

\else
    % =============================================================================
% AGENTE  : Geneticista (Especialista em Domínio) — AUTOR PRINCIPAL DESTE MÓDULO
% MÓDULO  : Discussão — Português Brasileiro (discussion_pt.tex)
% NOTA    : Tradução de discussion_en.tex — o conteúdo técnico é de
%           responsabilidade do Geneticista; a tradução é do Cientista-Líder.
% STATUS  : RASCUNHO
% =============================================================================

\section{Discussão}
\label{sec:discussion}

Os resultados obtidos com o painel simulado de marcadores estão em boa
concordância com as expectativas estabelecidas na literatura de marcadores
moleculares. Os valores consistentemente mais elevados de \pic{} e \he{}
observados para os marcadores \ssr{} corroboram sua reconhecida superioridade
na captura de polimorfismo, característica atribuível ao modelo de mutação
em degraus que rege a evolução de repetições de microssatélites e ao
consequente acúmulo de múltiplos alelos por locus \autocite{Ellegren2004}.

Os marcadores \snp{}, embora individualmente menos informativos devido à sua
natureza bialélica, compensam por sua abundância no genoma. Plataformas de
genotipagem de alta densidade de \snp{} podem capturar padrões de
desequilíbrio de ligação em alta resolução, viabilizando estudos de
associação genômica ampla (GWAS) e seleção genômica com acurácia superior
à alcançável com painéis mais esparsos de \ssr{} \autocite{Goddard2007}. O
equilíbrio entre a informatividade por locus (\ssr{}) e a densidade de
marcadores (\snp{}) deve, portanto, ser avaliado conforme a aplicação
específica: análise de parentesco e identificação de cultivares tipicamente
favorecem o \ssr{}, enquanto predição genômica e mapeamento de QTLs em
escala favorecem matrizes densas de \snp{}.

A distribuição estreita dos valores de $H_o$ em torno de $H_e$ no conjunto
de dados simulado indica uma população em equilíbrio de Hardy--Weinberg
aproximado, conforme esperado dada a ausência de seleção, deriva ou
endogamia no modelo de simulação. Em populações reais, desvios desse
equilíbrio podem sinalizar fenômenos biológicos de interesse, como depressão
por endogamia, varreduras seletivas ou artefatos de genotipagem (por exemplo,
alelos nulos em locos \ssr{}), os quais devem ser sistematicamente testados
antes das análises subsequentes \autocite{Peakall2006}.

Em conjunto, esses resultados reforçam a complementaridade dos sistemas de
marcadores \snp{} e \ssr{}. O pipeline analítico modular desenvolvido aqui,
combinando computação estatística em Python com geração de documentos
\LaTeX{}, fornece um arcabouço reprodutível e extensível para estudos
futuros de caracterização de marcadores.

\subsection{Limitações e Perspectivas Futuras}

A presente análise é limitada pelo uso de dados simulados e por um modelo
bialélico para os locos \ssr{}, que subestima seu verdadeiro polimorfismo.
Trabalhos futuros devem incorporar dados de genotipagem empírica de
populações reais, empregar formulações multialélicas de \pic{} para
marcadores \ssr{} \autocite{Anderson1993}, e estender o arcabouço
estatístico para incluir estatísticas $F$ e análises de estrutura
populacional (\textit{e.g.}, STRUCTURE, ADMIXTURE).

\fi

% --- 5. CONCLUSION ---
\ifnum\pdfstrcmp{\langversion}{en}=0
    % =============================================================================
% AGENT   : Lead Scientist
% MODULE  : Conclusion — English
% =============================================================================

\section{Conclusion}
\label{sec:conclusion}

This study demonstrates a reproducible, multi-agent pipeline for the
statistical characterisation and scientific reporting of molecular marker
panels. The simulated dataset confirmed that \ssr{} markers display
higher per-locus \pic{} and expected heterozygosity values than \snp{}
markers, while the latter offer advantages in marker density and
genotyping scalability.

The modular \LaTeX{} document structure, combined with automated Python
analysis scripts, enables rapid iteration as new data become available
or journal formatting requirements change. This framework will be applied
to empirical genotyping data from [target species] in subsequent phases
of this research programme.

\else
    % =============================================================================
% AGENTE  : Cientista-Líder
% MÓDULO  : Conclusão — Português Brasileiro (conclusion_pt.tex)
% SINCRON.: Deve ser consistente com conclusion_en.tex
% STATUS  : RASCUNHO
% =============================================================================

\section{Conclusão}
\label{sec:conclusion}

Este estudo demonstra um pipeline reprodutível e multi-agente para a
caracterização estatística e o relato científico de painéis de marcadores
moleculares. O conjunto de dados simulado confirmou que os marcadores \ssr{}
apresentam valores de \pic{} e heterozigose esperada superiores por locus
em comparação aos marcadores \snp{}, enquanto estes últimos oferecem
vantagens em termos de densidade de marcadores e escalabilidade da
genotipagem.

A estrutura modular do documento \LaTeX{}, combinada com scripts de análise
Python automatizados, permite iteração rápida à medida que novos dados se
tornem disponíveis ou que os requisitos de formatação do journal mudem.
Este arcabouço será aplicado a dados de genotipagem empírica de [espécie
alvo] nas fases subsequentes deste programa de pesquisa.

\fi

% ---------------------------------------------------------------------------
% ACKNOWLEDGEMENTS
% ---------------------------------------------------------------------------
\ifnum\pdfstrcmp{\langversion}{en}=0
    \section*{Acknowledgements}
    The authors thank [Funding Agency] (Grant No. XXXX) for financial
    support, and the Bioinformatics Laboratory of [Institution] for
    computational infrastructure.
\else
    \section*{Agradecimentos}
    Os autores agradecem à [Agência de Fomento] (Processo n.º XXXX) pelo
    apoio financeiro e ao Laboratório de Bioinformática de [Instituição]
    pela infraestrutura computacional.
\fi

% ---------------------------------------------------------------------------
% REFERENCES
% ---------------------------------------------------------------------------
\ifnum\pdfstrcmp{\langversion}{en}=0
    \printbibliography[heading=bibintoc, title={References}]
\else
    \printbibliography[heading=bibintoc, title={Referências}]
\fi

% ---------------------------------------------------------------------------
% SUPPLEMENTARY MATERIAL
% ---------------------------------------------------------------------------
\clearpage
\appendix

\ifnum\pdfstrcmp{\langversion}{en}=0
    \section{Supplementary Figures}
\else
    \section{Figuras Suplementares}
\fi

\begin{figure}[H]
    \centering
    \includegraphics[width=0.85\textwidth]{../output/figures/test_plot.pdf}
    \caption{%
        \ifnum\pdfstrcmp{\langversion}{en}=0
            Boxplot of Polymorphism Information Content (\pic{}) and expected
            heterozygosity (\he{}) for simulated \snp{} and \ssr{} markers
            (\textit{n}~=~60). Horizontal lines indicate medians; boxes
            represent the interquartile range (IQR); whiskers extend to
            1.5\,×\,IQR; circles denote outliers. $\mu$ = group mean.
        \else
            Boxplot do Conteúdo de Informação sobre Polimorfismo (\pic{}) e
            heterozigose esperada (\he{}) para marcadores \snp{} e \ssr{}
            simulados (\textit{n}~=~60). As linhas horizontais indicam as
            medianas; as caixas representam o intervalo interquartílico (IQR);
            os bigodes se estendem a 1{,}5\,×\,IQR; círculos indicam
            \textit{outliers}. $\mu$ = média do grupo.
        \fi
    }
    \label{fig:pic_boxplot}
\end{figure}

\clearpage

\ifnum\pdfstrcmp{\langversion}{en}=0
    \section{Supplementary Tables}
\else
    \section{Tabelas Suplementares}
\fi

\input{../output/tables/marker_summary}

\end{document}
